%% =====================================================================
%%  8 Strengths and Weaknesses
%% =====================================================================
%%  GUIDE
%%   - Required by the rules: discuss the strengths and weaknesses of your
%%     work. 38 of 74 gold papers do it explicitly; most others bury it.
%%   - Not "the model has errors", but "neglecting X biases Y by Z%".
%%   - Tie each item to a model component, and give at least one weakness
%%     with a number attached -- it reads as competence, not confession.
%% =====================================================================
\section{Strengths and Weaknesses}
\label{sec:sw}

\subsection*{Strengths}

\begin{enumerate}[label=(\arabic*),leftmargin=*]
  \item The model is built in three explicit layers, and each layer is
        priced: wall heat transfer moves $t_{0.3}$ by \SI{2}{\percent} at
        \SI{10}{cm} and \SI{5}{\percent} at \SI{5}{cm}, so a reader can see
        what was gained rather than take the complexity on trust.
  \item Every assumption carries a quantitative validity check
        (Table~\ref{tab:assumptions}); none is merely asserted, and the
        one that is weakest ($C_d$ constancy) is the one that is
        calibrated.
  \item The result is reduced to a scaling law and a single dimensionless
        time, $t_{0.3}/\tau=1.41$, so the crew can evaluate a different
        module without running the code.
  \item The numerical error is two orders of magnitude below the physical
        error (\SI{0.09}{\percent} against \SI{12}{\percent}), which is
        where it should be for a problem of this kind.
\end{enumerate}

\subsection*{Weaknesses}

\begin{enumerate}[label=(\arabic*),leftmargin=*]
  \item The wall heat transfer is first order with a single lumped time
        constant; with $\tau_w$ uncertain by \SI{30}{\percent} the answer
        moves by under \SI{1}{\percent} at both \SI{10}{cm} and
        \SI{5}{cm}, so the thermal layer earns its place through the
        trend in the exponent rather than through the nominal number.
  \item The lumped $C_d$ absorbs the throat density ratio and the
        contraction, so its fitted value cannot be compared with handbook
        data; the model is therefore calibrated for a sharp-edged circular
        hole and would need a new fit for a torn or irregular breach.
  \item The single-volume assumption fails for a compartmented module,
        so the model must not be applied there without a network
        extension, and it says nothing about the direction of the jet or
        the forces on the structure once the flow is established.
\end{enumerate}
